% Verbatim from minor_report/chapters/chapter_3_methodology.tex lines 303-530 (retired report),
% headings demoted by 0 level(s). Do not edit here; edit the source of the numbers.

Every simulation result in this report rests on how faithfully the simulator
reproduces the sensors and the contact physics. This subsection states what is
modelled, what is idealised, and in which direction each idealisation biases the
comparison between the visual and proprioceptive estimators.

\paragraph{Inertial measurement.} \Cref{tab:imu-noise-assumed} reports the noise
measured from the recorded stream and queries it, because no
\repo{dataset_allsensor} bag contains a stationary interval long enough for a
sound estimate. That query can be closed without a new recording: the noise is
\emph{declared}, not emergent, so the simulator's own configuration is the
ground truth the measurement was trying to recover.

\Cref{tab:injected-noise} collects what the simulator actually injects, together
with the measured values and the EuRoC reference. The measured gyroscope range
brackets the declared value, which corroborates both. The measured accelerometer
range does not: it is three to seven times the declared figure, and that
discrepancy is the estimator's rather than the sensor's, because the residual on
the driving axis carries real longitudinal acceleration that a
\SI{1.54}{\second} window cannot separate from noise. Prefer the declared values
wherever a number is needed.

\begin{table}[htbp]
\centering
\caption[Noise injected by the simulator, from the \repo{<imu>} and \repo{<camera>} sensor blocks of the robot model.]{Noise injected by the simulator, from the \repo{<imu>} and
\repo{<camera>} sensor blocks of the robot model. Densities are per-sample
$\sigma$ divided by $\sqrt{f_s}$ at \SI{200}{\hertz}, in the channel's own units
per $\sqrt{\si{\hertz}}$. The EuRoC column is the ADIS16448 of that
benchmark, on which both visual estimators were published.}
\label{tab:injected-noise}
\small
\setlength{\tabcolsep}{4pt}
\begin{tabularx}{\textwidth}{>{\raggedright\arraybackslash}X llll}
\toprule
Channel (units of $\sigma$) & Declared $\sigma$ & Density & Measured & EuRoC \\
\midrule
Gyroscope, per axis (\si{\radian\per\second}) & \num{0.0048} & \num{3.394e-4} &
  \numrange{3.00e-4}{3.64e-4} & \num{1.70e-4} \\
Accelerometer, per axis (\si{\metre\per\second\squared}) & \num{0.0566} &
  \num{4.002e-3} & \flag{\numrange{1.39e-2}{2.94e-2}} & \num{2.00e-3} \\
Camera, per pixel (normalised intensity) & \num{0.005} & --- & not measured & --- \\
\midrule
\multicolumn{5}{l}{\emph{Bias, first-order Gauss--Markov: $\sigma_b$ and
  correlation time $\tau$}} \\
Gyroscope bias (\si{\radian\per\second}) & \num{3.8785e-5} &
  $\tau=\SI{1000}{\second}$ & not measurable & --- \\
Accelerometer bias (\si{\metre\per\second\squared}) & \num{0.006} &
  $\tau=\SI{300}{\second}$ & not measurable & --- \\
\bottomrule
\end{tabularx}
\end{table}

Two readings follow. The simulated IMU is \emph{twice as noisy} as the EuRoC
reference on both channels, so it is not an optimistic sensor. The camera, at
about \num{1.3} grey levels out of \num{255} and with no distortion terms, is
very much cleaner than any real one. Both bias correlation times exceed the
\SIrange{85}{130}{\second} run lengths, so within a run the bias behaves as a
slowly varying constant and its instability is never exercised.

The IMU is therefore not optimistic in its noise. Three things are nonetheless
absent: chassis vibration, which on a real platform is frequently the dominant
error; scale-factor, cross-axis and $g$-sensitivity errors; and temperature
drift. The correlation times also exceed the \SIrange{85}{130}{\second} run
lengths, so within a run the bias behaves as a slowly varying constant and its
instability is never exercised. Finally, the message's \repo{orientation} field
carries no noise block at all and is therefore ground truth; no estimator in this
project consumes it, and none should.

\paragraph{Wheel encoders.} \repo{JointStatePublisher} reports the true solver
angle. Slip consequently appears as a \emph{discrepancy} rather than an error ---
the joint angle stays true, so encoder-integrated distance over-reads whenever
traction breaks, which is exactly a real encoder's behaviour. What is idealised
is resolution: a physical 1024\,CPR encoder quantises to \SI{6.1e-3}{\radian},
setting a floor on low-speed velocity estimation that the simulation does not
have. \repo{script/encoder_sim.py} exists to impose that floor downstream and has
not been applied to any result in this report.

\paragraph{Physics engine and solver.} \Cref{tab:physics-setup} records the
simulation's physics configuration, which is identical across all 26 world files.
Two entries deserve comment because they are easy to misread.

The \repo{<physics>} element declares \repo{type="ignored"}. That is not a
misconfiguration: in Gazebo Fortress the engine is selected by the physics system
plugin, not by that attribute, and the attribute is disregarded. No world passes
an \repo{<engine>} parameter to the plugin, so the engine is Fortress's default,
\repo{dartsim}. This matters for the friction discussion below, and it is the
reason \repo{slip1} and \repo{slip2} have no effect.

No solver block is declared anywhere --- no iteration count, no SOR parameter, no
constraint-force mixing or error-reduction terms. Every one of those therefore
takes the engine default. A \SI{1}{\milli\second} step is a fine integration
step for a \SI{1.4}{\metre\per\second} ground vehicle, but the untuned contact
solver is the part of this setup least able to support a quantitative claim about
contact behaviour.

\begin{table}[htbp]
\centering
\caption[Gazebo physics configuration, identical in all 26 world files.]{Gazebo physics configuration, identical in all 26 world files. Entries
marked ``engine default'' are not declared anywhere in the project and are
therefore whatever \repo{dartsim} chooses.}
\label{tab:physics-setup}
\small
\setlength{\tabcolsep}{4pt}
\begin{tabularx}{\textwidth}{>{\raggedright\arraybackslash}p{0.24\textwidth}>{\raggedright\arraybackslash}p{0.22\textwidth}>{\raggedright\arraybackslash}X}
\toprule
Setting & Value & Note \\
\midrule
Engine & \repo{dartsim} & Fortress default; no world sets \repo{<engine>} \\
\repo{<physics type>} & \repo{"ignored"} & disregarded in Fortress; the plugin
  selects the engine \\
\repo{max_step_size} & \SI{1}{\milli\second} & \SI{1}{\kilo\hertz} integration \\
\repo{real_time_factor} & 1 & no deliberate slow-down \\
\repo{real_time_update_rate} & engine default & not declared \\
Solver type and iterations & engine default & no \repo{<solver>} block anywhere \\
CFM / ERP & engine default & no \repo{<constraints>} block anywhere \\
Gravity & \SI{-9.8}{\metre\per\second\squared} & against ORB-SLAM3's hard-coded
  \num{9.81}, a \SI{0.1}{\percent} difference \\
\repo{slip1}, \repo{slip2} & 0 on all grounds & \flag{an ODE feature;
  \repo{dartsim} does not implement it, so inert at any value} \\
\bottomrule
\end{tabularx}
\end{table}

Achieved real-time factor was not logged during recording, so the table states
the requested factor rather than the realised one. Replay rate, which is logged,
is a separate control and is given per experiment in \cref{chap:results}.

\paragraph{Contact friction, and why the encoder is flattered.} The wheels
declare $\mu=1$, $\mu_2=1$ and every ground model declares $\mu=100$, $\mu_2=50$.
Both sets sit inside \repo{<friction><ode>} elements, which is where SDF places
them irrespective of engine, and \repo{dartsim} consumes them
(\cref{tab:physics-setup}). How it combines two differing coefficients is not
established here, so the effective value is not asserted; what is measured is the
outcome in \cref{tab:measured-slip}.

\repo{slip1} and \repo{slip2} are pinned at zero on all three ground models, but
that is not why force-dependent slip is absent. Those parameters are an ODE
feature, and \repo{dartsim} does not implement them, so they would be inert at
any value. Contact is rigid until Coulomb traction is exceeded regardless.

\Cref{tab:measured-slip} gives the consequence, measured rather than assumed, by
integrating rear-wheel angle against true path length.

\begin{table}[htbp]
\centering
\caption[Measured wheel slip on the five multi-sensor recordings.]{Measured wheel slip on the five multi-sensor recordings. The ratio is
encoder-integrated distance over true path length; above one is slip, below one
is a rolling-radius error and cannot be slip. Measured by
\repo{script/check_slip.py}.}
\label{tab:measured-slip}
\small
\begin{tabular}{lrrr}
\toprule
Recording & Encoder (m) & Truth (m) & Ratio \\
\midrule
\repo{allsensor_000} & 117.12 & 117.05 & 1.0006 \\
\repo{allsensor_001} & 112.91 & 113.11 & 0.9982 \\
\repo{allsensor_002} & 115.37 & 115.20 & 1.0015 \\
\repo{allsensor_003} & 141.68 & 140.61 & 1.0076 \\
\repo{allsensor_004} & 142.36 & 142.25 & 1.0008 \\
\bottomrule
\end{tabular}
\end{table}

Slip is present but under \SI{1}{\percent} everywhere, and on
\repo{allsensor_001} the encoder under-reads, which cannot be slip. Against a
wheel-odometry drift of \SIrange{2.15}{6.57}{\percent}, slip is therefore not the
dominant error; the yaw-rate geometry is, as the falsification of H1 establishes.

\begin{queried}[Threat to validity]
This is a bias with a known direction. The EKF baseline of
\cref{sec:results-proprio} is evaluated on a floor where the wheels effectively
cannot slip and with an encoder of unlimited resolution, which is dead
reckoning's best case. Both idealisations favour the EKF over the visual
estimators, and neither has been relaxed. The conclusion that a wheel+IMU filter
outperforms both visual pipelines should be read as holding \emph{under these
conditions}, not as a general result. \repo{script/make_slippery_floor.py}
generates ground models at a chosen friction coefficient so that traction can be
made an experimental factor; \repo{script/encoder_sim.py} imposes a realistic
quantisation floor. Neither has been used in any result reported here.
\end{queried}

\paragraph{Noise models supplied to the estimators.} The two are not told the
truth about the IMU, as \cref{tab:configured-vs-injected} shows.

\begin{table}[htbp]
\centering
\caption[What the visual estimators are configured to expect, against what the simulator injects.]{What the visual estimators are configured to expect, against what the
simulator injects. Both estimators use identical values. Inflating noise above
datasheet figures is standard practice in visual-inertial odometry and the
configuration records that intent; the ratios are given so the size of the
margin is visible.}
\label{tab:configured-vs-injected}
\small
\begin{tabularx}{\textwidth}{llll>{\raggedright\arraybackslash}X}
\toprule
Parameter & Configured & Injected & Ratio & Quantity \\
\midrule
\repo{gyr_n} & \num{0.002}  & \num{3.394e-4} & \num{5.9}$\times$ &
  gyroscope noise density \\
\repo{acc_n} & \num{0.02}   & \num{4.002e-3} & \num{5.0}$\times$ &
  accelerometer noise density \\
\repo{gyr_w} & \num{0.0001} & \num{1.735e-6} & \flag{\num{58}$\times$} &
  gyroscope bias random walk \\
\repo{acc_w} & \num{0.001}  & \num{4.899e-4} & \num{2.0}$\times$ &
  accelerometer bias random walk \\
\bottomrule
\end{tabularx}
\end{table}

The injected bias random walks are the Gauss--Markov driving densities
$\sigma_b\sqrt{2/\tau}$ implied by \cref{tab:injected-noise}. Inflating noise above datasheet values is
standard practice in visual-inertial odometry, and the configuration records that
intent. The EKF baseline, by contrast, derives its noise model by measuring it
against ground truth from the same recording.

That asymmetry was tested directly rather than left as a caveat. Configurations
carrying the true injected densities --- \repo{stereo_imu_truenoise.yaml} for
both estimators --- were run on all five recordings on 18 September 2026, after
the reporting cutoff. Accuracy improved wherever a run completed: VINS-Fusion on
four of five, by \SI{43}{\percent} and \SI{32}{\percent} on two of them. But
three of ten runs diverged outright, reaching \SIrange{111}{3899}{\metre} of
error with \numrange{516}{1545} tracking jumps, where the inflated configuration
diverged in none. The behaviour is bimodal rather than graded: surviving runs
carry zero to three jumps, failed runs hundreds. This confirms the configuration
comment's own claim that too-large noise degrades gracefully while too-small does
not, and the inflated setting is retained for every result in this report. These
runs are recorded as \repo{logging_20260916_03} and are not part of the reported
campaign.
